%!TEX root = ../main.tex
% =============================================================================
%  二、模型假设与符号说明        建议 2 页
%  评审在问：你的假设合理吗？符号清楚吗？
%  本节必须出现：① 假设表 A1–A10，每条含「必要性 + 现实/题目依据 + 影响方向 + 验证方式」
%                ② 信息结构 R1/R2/R3 的显式声明与实测对照
%                ③ 完整符号表，含单位，且必须含对偶变量
%                ④ 时标与单位口径声明（位置读法 / 电量口径 / Δt / E^max）
%  判据：一条假设若不能回答“不用它模型会怎样”，它就不该出现在假设表里。
% =============================================================================

\section{模型假设与符号说明}

\subsection{模型假设}

题面已给定储能最大容量 \(12000\)~kWh、最大充放电功率 \(5000\)~kW、2025 年 1 月 1 日
\(0{:}00\) 储电量 \(6000\)~kWh、储电量运行区间 \(1200\)--\(10800\)~kWh 与充放电效率
\(90\)\%；本文取运行区间为额定容量的 \(10\)\%--\(90\)\%（可用摆幅
\(10800-1200=9600\)~kWh），二者与功率上限构成两类硬约束，谁更紧由第四章容量/功率双向扫描
判定。

\begin{table}[htbp]
  \centering
  \caption{模型假设及其必要性、现实/题目依据、影响方向与验证方式（A1--A10）}
  \label{tab:assumptions}
  \footnotesize
  \setlength{\tabcolsep}{3pt}
  \begin{tabularx}{\linewidth}{C{0.65cm} L{2.45cm} L{2.95cm} L{2.6cm} L{1.55cm} Y}
    \toprule
    编号 & 假设 & 必要性（不用会怎样） & 现实/题目依据 & 影响方向 & 验证或灵敏度方式 \\
    \midrule
    A1 & 不允许向外部电网售电
       & 目标函数将多出一项售电收益，盈余有了出路，弃光量恒为 0，模型的可行域与最优值都改变
       & 题面未给出上网电价；就近消纳“优先自用”的政策导向
       & 使费用\textbf{偏高}，结果保守
       & 对结果文件作非负性审计（审计项 A7）：五个结果文件的负购电量数均为 0；
         并对引入上网电价 \(\theta p_t\) 的政策情景作扫描 \\
        A2 & 充电效率与放电效率\textbf{单边各取} \(0.9\)（往返 \(0.81\)）
       & 效率决定套利是否有利可图，是收益的核心参数；不声明读法则同一模型可给出三套结果
       & 题面给出“充放电效率为 \(90\%\)”，\textbf{未区分充、放方向}
       & 效率越低，往返损耗越大，储能套利收益越低
       & 三种主口径读法的节省额对照（单边各 \(0.9\)、仅充电侧 \(0.9\)、仅放电侧 \(0.9\)）
         见本节末段；另报“往返 \(0.9\) 等效单边”口径作条件对照 \\
        A3 & 同一时段不同时充放电
       & 若允许同时充放，须引入 \(0\)--\(1\) 变量，模型由线性规划变为混合整数规划，
         规模与求解时间上升，且对偶信息不再直接可用
       & 效率乘积 \(\eta_c\eta_d<1\) 时同时充放严格劣于只做净动作（命题~\ref{prop:no-simul}）
       & 不影响最优目标值，只影响求解方式
       & 数值断言 \(\min(c_t,d_t)=0\) 全部成立：问题 1 互斥违反数为 0，
         问题 2/3/4 的内部解互斥违反数也均为 0（审计项 A5） \\
        A4 & 忽略自放电与循环衰减
       & 否则目标函数须加入与循环次数（放电量）相关的衰减成本项，并需要衰减参数
       & 题面未给出任何衰减参数；结果窗口仅 334 天
       & 高估长期收益，且放电量越大高估越多
       & 加入循环成本 \(c_{cyc}\)（元/kWh 放电量）作单参数灵敏度，给出“储能节省额归零”的
         临界 \(c_{cyc}\) \\
        A5 & 储电量\textbf{跨日连续递推}，不逐日清零
       & 若每日清零，则与“2025-01-01 \(0{:}00\) 储电量 \(6000\)~kWh”的初值条件矛盾，
         并丢弃跨日套利空间
       & 附录 1 给出的初值 \(+\) 储能设备连续运行的常识；题面只要求 \(0{:}00\) 与
         \(24{:}00\) 储电量相同
       & 影响全窗口费用与每天的储电量初值，跨日套利越大影响越大
       & 与“每日重置初值”变体对照（消融层级 B2/B3）；审计项 A6 逐日核对
         “末值不低于初值”的违反天数（问题 2/3/4 均为 0） \\
        A6 & 每日终端条件 \(S_{144}\ge S_0\)；问题 1 取严格日周期 \(S_0=S_{144}=6000\)~kWh
       & 不设终端条件时，模型会在每天末尾与窗口末尾“掏空”储能以压低当期费用，
         形成短视解，且该解在真实连续运行中不可执行
       & 储能运营常识；问题 1 题面直接要求 \(S_0=S_{144}\)
       & 影响每天末尾与窗口末尾的费用；本数据上日周期口径比自由环状口径贵
         \(8.35\) 元（\(0.02\%\)）
       & 终端处理四组对照：完全自由、\(S_T\ge S_0\)、日周期、水价值终端（§9.4 层级 B5） \\
        A7 & 调整决策“后覆盖前”，调整费用的计价基准始终为 \(0{:}00\) 的\textbf{计划购电量}
       & 若基准随上一次调整漂移，调整费与违约费的口径不可复现，也无法与题面表述对应
       & 题面“计划购电量高于调整购电量的部分”一句明确了基准
       & 影响调整与违约相关费用；四次预报口径下违约费 \(210\,674.71\) 元、
         超额费 \(262\,628.79\) 元
       & 与“相对上一次调整量”口径并列报告并各自命名（第九章） \\
        A8 & 2025 年 1 月作为储电量递推的\textbf{预热期}，其终值即 2025-02-01 \(0{:}00\) 的初值
       & 结果窗口自 2025-02-01 起，而初值自 2025-01-01 起算；若不作预热声明，
         2 月 1 日的储电量没有来源，跨日递推无从开始
       & 附件 2 覆盖 2025-01-01 至 12-31，结果文件要求 2025-02-01 起共 334 天
       & 影响窗口初期的储电量与初期费用；1 月本身不计入任何结果
       & 论文与结果文件中明示 1 月不计入结果，并公开 2 月 1 日初值的来源；
         审计项 A6 逐日核对终端条件 \\
        A9 & \textbf{信息结构取 R3}：问题 2/4-2 的 \(0{:}00\) 计划依据为附件 3 中当日 \(0{:}00\)
         发布的光伏预报，负载曲线按“已知”处理；问题 3/4-3 另有 \(6{:}00/12{:}00/18{:}00\)
         三次滚动的当次预报
       & 题面只说“每天 \(0{:}00\) 制定计划”，未说明 \(0{:}00\) 知道什么；不声明则
         问题 2--4 的任何数值都无法解释（三种解读的窗口总费用相差近千万元）
       & 题面“根据预测数据制定计划”；附件 3 提供 \(0{:}00\) 发布的 24~h 光伏预报；
         题面未给负载预报，故按已知处理
       & 费用落在完美信息（下界）与典型日预报（上界）之间；负载按已知处理会
         \textbf{低估}紧急购电量，使结果偏乐观
       & 三种信息结构 R1/R2/R3 在同一 334 天窗口与同一场景集上的完整对照
         （表~\ref{tab:infostruct}）；\textbf{非预期性置换检验}：置换 \(6{:}00/12{:}00/18{:}00\)
         发布的预报后，其发布时刻之前各槽位的计划量逐位不变（最大差值 \(0.0\)），
         置换 \(0{:}00\) 预报则只改变 \(0{:}00\) 之后的决策（第九章）；
         并对预报误差水平另作灵敏度 \\
        A10 & 波动电价（问题 4）下 \(0{:}00\) 可用的\textbf{电价预报取“分时骨架”}，
         即电价的逐时段日内均值曲线；结算仍按附件 4 的实际电价
       & 附件 4 未给出日前电价预报；若不声明电价预报口径，问题 4 的“计划”将使用
         结算时的实际电价，构成数据泄漏，结果会被系统性压低
       & 分时骨架与附件 4 的逐时段均值最大偏差仅 \(5.15\times10^{-5}\)~元/kWh，
         且等于附件 1 的电价剖面，是可获得的日前信息
       & 忽略电价的残差波动，使问题 4 的费用与节省率\textbf{偏乐观}
       & 结算侧改用附件 4 实际电价，价格随机性只在结算体现；并以“是否保留残差自相关”
         作对照（分块 Bootstrap 条件对照的分位数最大相对偏差达 \(20.44\%\)） \\
    \bottomrule
  \end{tabularx}
\end{table}

假设 A2 影响最大：三种效率读法下，典型日相对无储能基线的节省额为
\(12\,925.10\)~元（单边各 \(0.9\)，\textbf{主口径}）、\(14\,635.52\)~元（仅充电侧 \(0.9\)）、
\(13\,883.44\)~元（仅放电侧 \(0.9\)），极差 \(13.2\%\)；把题面“\(90\%\)”读作往返效率、等效
为单边 \(\sqrt{0.9}\) 时为 \(14\,250.99\)~元。主口径写入全部结果表，其余读法只在灵敏度节出现
（\evid{C-Q1-03}）。

\subsection{信息结构声明}

问题 2--4 中“每天 \(0{:}00\) 制定计划”有三种解读，本文各跑一次全窗口实验，条件完全一致
（同一 334 天窗口、同一场景集、同一终端条件、同一电价口径），只有“\(0{:}00\) 已知什么”
不同，三行可直接相减（表~\ref{tab:infostruct}）。

\begin{table}[htbp]
  \centering
  \caption{信息结构的三种解读及其在 2025-02-01 至 12-31（334 天）上的实测后果
           （费用单位：元，电量单位：kWh，电价均为附件 1 固定价；
           紧急购电量口径见 \calref{C-05}，“出现在多少天”一并列出）}
  \label{tab:infostruct}
  \footnotesize
  \setlength{\tabcolsep}{3pt}
  \begin{tabularx}{\linewidth}{C{0.8cm} L{3.3cm} C{2.55cm} C{1.5cm} C{2.3cm} Y}
    \toprule
    代号 & 解读 & 总费用 & 相对基线节省率 & 紧急购电量 & 本文的使用方式 \\
    \midrule
    R1 & 完美信息：\(0{:}00\) 即知当日负载与光伏真值
       & \(13\,411\,560.90\) & \(18.26\%\) & \(0.00\)（0/334 天）
       & \textbf{不可作为主方案}；用作“预测误差的代价”的\textbf{下界对照} \\
        R2 & 典型日预报：以附件 1 的日均形状作为预报
       & \(23\,484\,646.63\) & \(-43.14\%\) & \(2\,525\,548.76\)（286/334 天）
       & 仅作悲观对照；附件 1 是实际值的全年逐时段均值而\textbf{不是预报}，
         用它当预报须额外论证并做泄漏检验，本文\textbf{不采用} \\
        R3 & 日前预报：\(0{:}00\) 使用附件 3 当日 \(0{:}00\) 发布的光伏预报，
         负载曲线按已知处理
       & \(14\,264\,792.13\) & \(13.06\%\) & \(172\,015.82\)（296/334 天）
       & \textbf{本文主方案}：最符合工程实际与题面“制定计划”的语义 \\
    \bottomrule
  \end{tabularx}
\end{table}

本文采用\textbf{信息结构 R3}：R1 假设 \(0{:}00\) 即知当天真值，紧急购电量恒为 0、问题 2
退化为问题 1；R2 用事后平均值冒充事前信息，总费用比无储能基线还高 \(43.14\%\)，
属口径错误；R3 用附件 3 中\textbf{当天 \(0{:}00\) 发布}的真实预报，全部信息先于决策时刻
可得，无泄漏也不夸大。三点补充：① 计划依据（附件 3 预报 \(+\) 负载曲线）与结算依据
（附件 2 当日实际负载与光伏）是\textbf{两套数据}，差额即紧急购电来源；② “预测误差的代价”
为 \(14\,264\,792.13-13\,411\,560.90=853\,231.23\) 元、占 R1 费用的 \(6.36\%\)
（\evid{C-Q2-01}）；③ 负载按“已知”处理是对 R3 的乐观简化（题面未给负载预报），会低估
紧急购电量，第九章另作预报误差灵敏度。

\subsection{符号说明}

公式符号均在表~\ref{tab:notation} 登记，单位口径见本节末段。

\begin{table}[htbp]
  \centering
  \caption{符号说明（集合、参数、决策变量、对偶变量与指标量；“---”表示无量纲）}
  \label{tab:notation}
  \footnotesize
  \setlength{\tabcolsep}{3pt}
  \begin{tabularx}{\linewidth}{C{1.2cm} C{2.55cm} L{5.1cm} C{1.35cm} Y}
    \toprule
    类别 & 符号 & 含义 & 单位 & 来源/取值 \\
    \midrule
    集合 & \(t\in\mathcal T\) & 时段（10~min 粒度） & --- & \(|\mathcal T|=144\) \\
    集合 & \(d\in\mathcal D\) & 日期 & --- & 334 天（2025-02-01 至 12-31） \\
    集合 & \(m\in\mathcal M\) & 预报发布/决策时刻 & --- & \(\{0,6,12,18\}\) 时 \\
    集合 & \(\omega\in\Omega\) & 情景 & --- & \(|\Omega|=K\) \\
        参数 & \(L_{d,t}\) & 小区负载（功率） & kW & 附件 2 \\
    参数 & \(PV_{d,t}\) & 光伏实际功率 & kW & 附件 2 \\
    参数 & \(\widehat{PV}^{(m)}_{d,h}\) & \(m\) 时发布的光伏预报功率（整点值） & kW & 附件 3 \\
    参数 & \(p_{d,t}\) & 电价 & 元/kWh & 附件 1（固定价）/ 附件 4（波动价） \\
    参数 & \(\bar p_t\) & 电价的日内分时骨架（逐时段均值） & 元/kWh & 附件 1 电价剖面（假设 A10） \\
    参数 & \(\varepsilon_{d,t}\) & 电价相对骨架的残差 & 元/kWh & \(p_{d,t}-\bar p_t\) \\
    参数 & \(\Delta t\) & 时间步长 & h & \(1/6\) \\
    参数 & \(\eta_c,\eta_d\) & 充电效率、放电效率（单边） & --- & \(0.9/0.9\)（假设 A2） \\
    参数 & \(\underline S,\overline S\) & 储电量运行下限、上限 & kWh & \(1200/10800\) \\
    参数 & \(\overline S^{\max}\) & 储能设备最大容量 & kWh & \(12000\) \\
    参数 & \(\overline P\) & 最大充放电功率 & kW & \(5000\) \\
    参数 & \(E^{\max}\) & 单槽充/放电量上限 & kWh & \(\overline P\Delta t=833.3333\) \\
    参数 & \(S^{init}\) & 初始储电量 & kWh & \(6000\) @ 2025-01-01 \(0{:}00\) \\
    参数 & \(\kappa\) & 紧急购电电价倍数 & --- & \(5\) \\
    参数 & \(\kappa^{+},\kappa^{-}\) & 调整超额 / 违约电价倍数 & --- & \(1.5/0.5\) \\
    参数 & \(\beta\) & CVaR 置信水平（预先锁定） & --- & \(0.95\) \\
    参数 & \(\pi_\omega\) & 情景权重 & --- & 分层样本占比，\(\sum_\omega\pi_\omega=1\) \\
    参数 & \(\rho(\cdot,\cdot)\) & 相关系数 & --- & 用于刻画电价与净负荷的共动 \\
        决策 & \(E^{plan}_{d,t}\) & 计划购电量（第一阶段，与 \(\omega\) 无关） & kWh & \(\ge0\) \\
    决策 & \(E^{adj}_{d,t}\) & 调整购电量（最后覆盖该槽位的决策） & kWh & \(\ge0\) \\
    决策 & \(g_{d,t}\) & 紧急购电量（第二阶段追索） & kWh & \(\ge0\) \\
    决策 & \(c_t,d_t\) & 充电量、放电量 & kWh & \(0\le c_t,d_t\le E^{\max}\) \\
    决策 & \(s_t\) & 弃光/盈余处置电量 & kWh & \(\ge0\) \\
    决策 & \(S_t\) & 时段末储电量 & kWh & \(\in[\underline S,\overline S]\) \\
    决策 & \(u_t,v_t\) & 调整量相对计划量的增、减分解 & kWh & \(E^{adj}_t-E^{plan}_t=u_t-v_t\)，均 \(\ge0\) \\
        \textbf{对偶} & \(\alpha_t\) & 电量平衡约束的影子价格（\(E_t>0\) 时等于 \(p_t\)） & 元/kWh & LP 对偶；实测取值 \([0.3713,1.2561]\)，恒等式偏差 \(0.000\) \\
    \textbf{对偶} & \(\nu_t\) & 储电量递推约束的乘子，即储电的边际价值（水价值） & 元/kWh & LP 对偶；实测 \(\nu_t\in[0.4711,1.1305]\) \\
    \textbf{对偶} & \(\mu_t\) & 水价值的等价记法（\(\mu_t\equiv\nu_t\)），第四章阈值律沿用 & 元/kWh & LP 对偶 \\
    \textbf{对偶} & \(\lambda^{\pm}_t\) & 储电量上/下界的影子价格 & 元/kWh & LP 对偶 \\
    \textbf{对偶} & \(\nu^{c}_t,\nu^{d}_t\) & 充/放电功率界的乘子 & 元/kWh & LP 对偶 \\
    \textbf{对偶} & \(z^{c}_t,z^{d}_t\) & 充/放互补松弛量 & 元/kWh & \(z^{c}_t=\alpha_t-\eta_c\nu_t\)；\(z^{d}_t=-\alpha_t+\nu_t/\eta_d\)（阈值律：\(\alpha_t\le\eta_c\nu_t\) 充电、\(\alpha_t\ge\nu_t/\eta_d\) 放电，144 槽违反次数 0） \\
        指标 & \(C^{base}\) & 无储能基线购电费 & 元 & 口径 \calref{C-01} \\
    指标 & \(C^{plan},C^{emg}\) & 计划购电费、紧急购电费 & 元 & 口径 \calref{C-03} \\
    指标 & \(V^{arb}\) & 储能与弃光带来的相对节省 & 元 & \(C^{base}-C^{*}\)，节省率按 \calref{C-04} \\
    \bottomrule
  \end{tabularx}
\end{table}

\paragraph{时标与单位口径声明}
本节符号遵循同一套时标与单位口径（四问一致）：① \textbf{位置读法}——附件 1/2/4 的标签为
\textbf{时段末端}，槽位 \(j\)（\(j=0,1,\dots,143\)）表示区间 \([10j,\ 10(j+1))\) 分钟，
故题面表 1 的“\(10{:}00\)--\(10{:}10\)”对应槽位 60；末列 \(0{:}00{+}1\) 属当日第 144 个槽位
\([23{:}50,24{:}00)\)、不得丢弃（逐格对照见 §3.2 表~\ref{tab:timeruler}）。② \textbf{电量
口径}——决策变量为电量（kWh），\(\Delta t=1/6\)~h，功率上限 \(\overline P=5000\)~kW 对应
单槽电量上限 \(E^{\max}=\overline P\Delta t=833.3333\)~kWh；目标函数的电价系数为
\(p_t\)（元/kWh）、\textbf{不再乘} \(\Delta t\)；等价的功率口径为“变量取功率、目标系数取
\(p_t\Delta t\)、上限取 \(5000\)~kW”，两套口径全文只用一套。③ \textbf{单位与精度}——功率
kW、电量 kWh、电价 元/kWh、费用 元；表格费用保留 2 位小数，内部计算不取整。
